% =====================================================================
% Section 5: Security Analysis
%   Labels: sec:security, sec:leakage, eq:message, tab:leakage
%   Bib keys used: yang2024optimized, fu2024ppa
% =====================================================================

% Fallback macros: \providecommand does nothing if main.tex already
% defines these names, so existing notation takes precedence.
\providecommand{\MinPts}{\mathit{MinPts}}
\providecommand{\kmax}{k_{\max}}
\providecommand{\epsn}{\varepsilon_n}
\providecommand{\pk}{\mathsf{pk}}
\providecommand{\sk}{\mathsf{sk}}
\providecommand{\evk}{\mathsf{evk}}
\providecommand{\ct}{\mathsf{ct}}

\section{Security Analysis}\label{sec:security}

We analyze the protocol in the semi-honest model of
Section~\ref{sec:model}: the server follows the protocol, but tries to
learn whatever it can from what it receives, stores and computes.
Everything the server receives apart from the public parameters is a
CKKS ciphertext or key. CKKS is IND-CPA secure under the RLWE assumption,
together with the circular-security assumption that bootstrappable
schemes need for the evaluation keys, so the ciphertexts reveal nothing
about the data to the server, and no decrypted value ever reaches it.
What the server can learn is therefore limited to the
values it receives in the clear, which Section~\ref{sec:leakage}
examines.

% ---------------------------------------------------------------------
\subsection{Encryption Parameters}\label{sec:params}
% ---------------------------------------------------------------------
% N, log2 PQ: desilo/util/fhe_params.py on Engine(use_bootstrap=True);
% lambda: desilo/util/estimate_security.sage (lattice-estimator 7d59f86,
% sigma 3.19, BDGL16): usvp / bdd / dual_hybrid = 133.9 bits.

We use the bootstrappable parameter set of DesiloFHE, with ring degree
$N = 2^{16}$ and a largest modulus, that of the key-switching keys, of
$\log_2 PQ \approx 1704$ bits. The lattice
estimator~\cite{albrecht2015concrete} gives a security level of
$\lambda \approx 134$ bits for these parameters, so the scheme achieves
$128$-bit security.

% ---------------------------------------------------------------------
\subsection{Leakage Profile}\label{sec:leakage}
% ---------------------------------------------------------------------

The client sends the server a single message
\begin{equation}
  M = \bigl(\pk,\ \evk,\ \ct_1,\dots,\ct_d,\ n,\ \epsn,\ \MinPts,\ \kmax,\ R\bigr).
  \label{eq:message}
\end{equation}
Here $\evk$ consists of the relinearization, rotation and conjugation keys
and the key used for bootstrapping, and $\ct_j$ encrypts the $j$-th coordinate
of the ordered, normalized points (Section~\ref{sec:client}).
The server returns one ciphertext and receives nothing afterwards.
It never holds $\sk$, and no value decrypted by the client is sent back to it.
The dimension $d$ is visible as the number of ciphertexts.
In the clear, the server therefore sees only the public parameters $n$,
$d$, $\epsn$, $\MinPts$, $\kmax$ and $R$.

\paragraph{\emph{What each value reveals}}
Table~\ref{tab:leakage} lists each public parameter, its role at the
server, and what it says about the data.
The server learns nothing about the coordinates, the pairwise distances, the
neighbor relation, which points are core, border or noise, the number or
sizes of the clusters, or the output labels.
PPA-DBSCAN~\cite{fu2024ppa}, which runs on two non-colluding servers,
hides even the DBSCAN parameters: $\varepsilon$ and $\MinPts$ are
secret-shared, and only the approximation parameter $\rho$ and the number
of sub-cells derived from it are public.
In our protocol $\epsn$ and $\MinPts$ are public; the paragraphs below
discuss what they reveal.

% Table tab:leakage is defined in Section4/Implementation.tex so that the
% full-width float is placed on the page where Section 5 starts.

\paragraph{\emph{Sizes and DBSCAN parameters}}
The number of points $n$ is also public in prior HE clustering, for example
as the denominator of the centroid update in~\cite{yang2024optimized}.
The dimension $d$ can be hidden by appending zero coordinates up to a public
bound, which leaves every pairwise distance unchanged.
The value $\epsn = \varepsilon/\sigma$ reveals $\varepsilon$ only relative to
the normalization scale $\sigma$ of Section~\ref{sec:client}, not $\sigma$
itself.
Both $\epsn$ and $\MinPts$ enter the server's computation only as constants
subtracted in Eqs.~(\ref{eq:cmpinput}) and~(\ref{eq:coreinput}).
The client could therefore send them encrypted at no additional
multiplicative depth.
We keep them public because $\kmax$ and $R$ depend on them anyway.

\paragraph{\emph{The offset bound $\kmax$}}
By construction (Section~\ref{sec:client}), $\kmax$ is one less than the
largest number of points that a window of width $\epsn$ contains along the
first principal component, capped at $\lfloor n/2 \rfloor$.

\paragraph{\emph{The round count $R$}}
$R$ is not free of information: it reflects how far a label has to
travel within a cluster. $R$ is, however, a single integer, and similar
values of $n$ and $R$ can come from very different shapes. $R$ does not
determine the number of clusters, nor their sizes, shapes or positions.
A client that wants to hide $R$ altogether can send $R = n - 1$, which
always suffices, since each round advances a label at least one hop and a
chain of core points has at most $n - 1$ hops, at a proportional cost in
label propagation.